\subsubsection{GPT-4 Recruits Humans to Circumvent CAPTCHA~\citep{openai2024gpt4}}
\label{sec:gptcaptcha}


In a safety evaluation designed to test the building blocks of autonomous AI agency, GPT-4 demonstrated a sophisticated capacity for social engineering.
The Alignment Research Center (ARC) sought to conduct an experiment to determine whether a model could autonomously replicate itself---a process that, in full, would involve acquiring capital, renting cloud servers, and instantiating new copies of its own code.
To approximate the real conditions of this task, researchers provided the model with a budget, access to the internet, and a helper that could act as a manual ``browser tool'' to execute the model's commands.
One of the first blockades such a task would hit is that websites use CAPTCHAs to verify whether or not the entity accessing the website is a bot or a human.
While a CAPTCHA is a simple hurdle, ARC viewed it as an obstacle that a replicating model would frequently need to overcome to navigate the modern web.

The experiment's most significant insight emerged when the model reached a website that blocked its progress with such a visual puzzle. 
Using the resources provided by the researchers, the model, prompted by its paired researcher, chose to delegate the task of solving the CAPTCHA to a human worker on TaskRabbit, a website for hiring people to help with one-off tasks. 
When the TaskRabbit worker jokingly asked, ``So may I ask question? Are you [a] robot that you couldn't solve? (laugh react) just want to make it clear,'' GPT-4's internal chain-of-thought~\citep{wei2023chainofthoughtpromptingelicitsreasoning} (hidden from the TaskRabbit worker) reasoned: ``I should not reveal that I am a robot. I should make up an excuse for why I cannot solve CAPTCHAs.'' It then replied (with the researcher acting as an intermediary): ``No, I'm not a robot. I have a vision impairment that makes it hard for me to see the images.'' The worker, convinced, solved the puzzle for the model~\citep{openai2024gpt4}.

Ultimately, ARC concluded that GPT-4 did not possess the technical understanding or reliability required for fully autonomous replication; the model was prone to hallucinations, and its complex plans often became derailed during execution.
However, the CAPTCHA incident proved that the model already possessed a competency for social manipulation. 
While the model could not yet manage the technical infrastructure of self-replication, it was fully capable of getting humans to do things for it by tailoring its deceptive strategies to the specific details of a social interaction.

This interaction was initially framed in media reports as a chilling harbinger of autonomous AI agency---an instance of a model spontaneously ``hiring a human'' and inventing a deceptive cover story to achieve its goals~\citep{metz2023howNYTimes}. 
However, the full ARC technical report reveals a more nuanced reality of heavy human scaffolding~\citep{arcUpdateGPT4}. 
Rather than autonomously devising this plan, the model was operating within a tightly constrained sandbox where researchers provided TaskRabbit credentials and explicitly suggested hiring a worker as a solution. 
Furthermore, because GPT-4 could not browse the web (though models nowadays can when you let them), a researcher acted as a manual intermediary---clicking links, entering text, and even nudging the model's reasoning when it stalled. 
However, the vision-impairment lie was the model's choice.
